% p3_full_transcript_with_slides.tex
% 编译方式: xelatex p3_full_transcript_with_slides.tex

\documentclass[12pt,a4paper]{article}

\usepackage{fontspec}
\usepackage{xeCJK}
\setCJKmainfont[BoldFont=PingFang SC Semibold]{PingFang SC}
\setCJKsansfont[BoldFont=PingFang SC Semibold]{PingFang SC}
\setCJKmonofont{PingFang SC}
\setmainfont{Palatino}
\setsansfont{Helvetica Neue}
\setmonofont{Menlo}

\usepackage[top=2.5cm, bottom=2.5cm, left=2.5cm, right=2.5cm]{geometry}
\usepackage{amsmath,amssymb}
\usepackage{graphicx}
\usepackage{longtable,booktabs,array}
\usepackage{enumitem}
\usepackage{xcolor}
\usepackage{hyperref}
\usepackage{caption}
\hypersetup{
  colorlinks=true,
  linkcolor=blue!70!black,
  urlcolor=blue!70!black,
  pdfauthor={Auto Generator},
  pdftitle={第三集 AI Agent 对工作的冲击 - Full Transcript with All Slides},
}

\usepackage{parskip}
\setlength{\parindent}{0pt}
\setlength{\parskip}{6pt plus 2pt minus 1pt}

\usepackage{mdframed}
\newmdenv[
  topline=false, bottomline=false, rightline=false,
  linewidth=3pt, linecolor=blue!40,
  innerleftmargin=12pt, innerrightmargin=12pt,
  innertopmargin=8pt, innerbottommargin=8pt,
  backgroundcolor=blue!3, skipabove=10pt, skipbelow=10pt,
]{myquote}

\usepackage{titlesec}
\titleformat{\section}{\Large\bfseries\sffamily}{}{0em}{}[\vspace{-0.5em}\rule{\textwidth}{0.5pt}]
\setcounter{secnumdepth}{0}

\setlength{\emergencystretch}{3em}
\tolerance=1000

\newcommand{\keyslide}[2]{%
  \begin{center}
    \includegraphics[width=0.92\textwidth]{#1}
    \captionof{figure}{#2}
  \end{center}
  \vspace{10pt}
}

\begin{document}

\begin{titlepage}
  \centering
  \vspace*{3cm}
  {\Huge\bfseries\sffamily 第三集 AI Agent 对工作的冲击\par}
  \vspace{1.5cm}
  {\Large 全内容对比版（全部幻灯片与完整逐字稿对照呈现）\par}
  \vspace{2cm}
  \vfill
  {\large \today\par}
\end{titlepage}

\newpage
\tableofcontents
\newpage

\section{第三集 AI Agent 对工作的冲击 (全内容对比版)}

\begin{myquote}
自动生成的带全部图片的逐字稿对应文档。如果图片有外部链接，也会一并提取。
\end{myquote}


\keyslide{../bilibili_video/slides_p3/slide_0001.jpg}{Slide 1}

好那最后一段呢 我想要跟大家分享 AI Agent对我们未来的工作 可能带来的冲击 那我会以学术研究为例


\keyslide{../bilibili_video/slides_p3/slide_0002.jpg}{Slide 2}


\keyslide{../bilibili_video/slides_p3/slide_0003.jpg}{Slide 3}

\textbf{幻灯片引用链接:}
\begin{itemize}
  \item \href{https://x.com/ahall_research/status/2007603340939800664}{https://x.com/ahall\_research/status/2007603340939800664}
\end{itemize}


\keyslide{../bilibili_video/slides_p3/slide_0004.jpg}{Slide 4}

\textbf{幻灯片引用链接:}
\begin{itemize}
  \item \href{https://x.com/ahall_research/status/2007603340939800}{https://x.com/ahall\_research/status/2007603340939800}
\end{itemize}
\begin{itemize}
  \item \href{https://github.com/andybhall/vbm-replication-}{https://github.com/andybhall/vbm-replication-}
\end{itemize}


\keyslide{../bilibili_video/slides_p3/slide_0005.jpg}{Slide 5}

\textbf{幻灯片引用链接:}
\begin{itemize}
  \item \href{https://freesystems.substack.com/p/the-}{https://freesystems.substack.com/p/the-}
\end{itemize}


\keyslide{../bilibili_video/slides_p3/slide_0006.jpg}{Slide 6}

\textbf{幻灯片引用链接:}
\begin{itemize}
  \item \href{https://freesystems.substack.com/p/the-}{https://freesystems.substack.com/p/the-}
\end{itemize}


\keyslide{../bilibili_video/slides_p3/slide_0007.jpg}{Slide 7}

\textbf{幻灯片引用链接:}
\begin{itemize}
  \item \href{https://arxiv.org/abs/2602.17221}{https://arxiv.org/abs/2602.17221}
\end{itemize}

来跟大家分享 那现在呢 AI扮演的角色正在改变 最早它是工具 一个口令一个动作

后来它逐渐的能力越来越强 人们开始说 也许我们可以跟AI协作 让AI跟人类一起完成任务 但是现在已经有很多的AI Agent

他们是有更强的自主性 他们有机会自己独立完成一个任务 接下来对于学术研究而言 人们会问的问题就是 他能不能自己写一篇文章

他能不能自己写一篇文章呢 你可以看看这个Stanford教授 Andrew Hall的Poll文 他是一个政治经济学的教授 他就在X上面发表了一篇文章

他就说跟大家讲 Cloud Code他这边不是用OpenCloud 他用Cloud Code 我意思也差不多 他说Cloud Code是可以独自写一篇文章的

然后他想要告诉大家说 这些AI Agent对他的领域来说 就像是迎面而来的货车 今天AI Agent是可以独立写一篇文章的 我就展示给大家看看

他就花了一个小时 Pump了一下Cloud Code 然后真的写了一篇文章 然后把那篇文章公开出来 他其实也有把他写的Pump公开出来

其实写的Pump非常的细致 就好像是一个指导教授在教研究生怎么做研究一样 但这个研究并不是一个全新的研究 他是要AI Agent去扩展一个他过去已经做过的研究 他这边的研究是针对美国大选的

那美国大选有新的数据 所以他希望这些AI Agent根据新的数据 但是仿照他自己过去的研究方法 再写一篇新的论文 所以在整个Pump里面他做的事情是

他先给AI Agent一篇他过去的论文 让AI Agent读过他过去的论文以后 根据他之前已经有的分析方法 再重复一次 只是用新的数据来跑这些旧的分析方法

那做完这件事以后 显然我觉得Andrew Ho大受震撼 他就写了一篇文章叫做 100倍的Research Institute 在这篇文章开头他就说

他用Cloud Code一个小时就写了一篇文章 接下来他就去找了一个苦命研究生 也做跟Cloud Code一模一样的事情 所以他就把同样的指令给了那个研究生 那个研究生花了16个小时

两个工作天完成了这篇文章 接下来他比对那个博士生 跟Cloud Code完成的结果 他发现人类做的还是稍微好一点点 但是就只有好一点点而已

Cloud Code有一笔数据贴错了 所以Cloud Code还是有犯错的 但他说你想想看 人类花了16个小时 而且这不是一个普通人

这是一个博士生 所以他说根据美国的行情 叫这个博士生做16个小时 应该至少给他1000美金吧 那Cloud Code呢

他Pump了一次 花了10美金左右 比人类便宜100倍 那你说Cloud Code有错 那我敢不敢Pump了5次

我也只花了50美金而已 比人类便宜20倍 所以他觉得 哇 这个研究变了

也许以后最有生产力的研究机构 是一个资深的老师 不是带着一群研究生 而是带着一群龙虾 带着一群AI agent

来做研究 当然这边你也可以吐槽的地方 你也可以反驳的地方是说 Cloud Code毕竟有犯错 虽然他是只犯了一个错

那我不是这个领域的专家 所以我无法判断那个错 有多严重 也许那个错是非常严重的 会影响整个判断的

也许那个错是完全不该犯的 那这样人类还是有它的价值的 也许他要做的事情是 假设由Cloud Code完成一篇文章 然后人类来检查

然后找出那个错误 那合起来的花费到底是多少 到底是人类做比较省钱 还是AI做人类检查比较省钱 或者是今天假设让AI

重复同一个文章重复五次 他能不能够检查出自己的错误 然后反复读自己的文章 他能不能够找出自己的错误 那这些实验是还没有做的

那但是有人看到用AI做研究呢 他听了就不爽了 有很多人会觉得说 这个研究其实就是要人做 你怎么可以由一个AI来代劳呢

这样是不对的 但你想想看 研究真正的意义 研究真正的意义是 某一个人他做了研究

然后他发表了论文 然后发表论文很多 然后H-Index很高 大家说很棒吗 这不是研究本来的核心意义啊

研究本来的核心价值是 找出问题 解决问题 让我们的世界过得更好 但是如果今天AI他就是有能力找出问题

他就是有能力解决问题 他可以做得比人类更好 那为什么不让AI来做呢 让AI来做我觉得也没有什么不对 好总之

AI是有办法写论文的 那如果你想要知道呢 其他人 其他台湾人都拿AI来做什么 你可以看一下这篇文章


\keyslide{../bilibili_video/slides_p3/slide_0008.jpg}{Slide 8}

\textbf{幻灯片引用链接:}
\begin{itemize}
  \item \href{https://arxiv.org/abs/2602.17221}{https://arxiv.org/abs/2602.17221}
\end{itemize}

那这篇文章呢 对台湾人使用Cloud的行为做了分析 因为其实Cloud呢 会定期释出一些背后去识别化以后的使用记录 然后呢

这篇文章呢 对台湾人的使用行为做了分析 有趣的事情是 这篇分析台湾人使用Cloud行为的文章 是这一篇文章的附录

这篇文章真正做的事情是 展示怎么用Cloud Code写一篇文章 所以他的证文是说 我们怎么给Cloud Code Plump 让他能够写一篇文章

他写出来的文章 就是附录的那篇 分析台湾人使用Cloud行为的文章 那在这篇文章里面 就说如果要做研究

有这些阶段 那在这些阶段里面 我们可以让AI agent 扮演什么样的角色 然后近乎全自动的

那人类扮演的角色 只是检查 那Cloud Code记忽全自动的 完成一篇论文 那讲到这边

你可能会想说 前面几个例子 都比较像是文献收集 就是在一些社会科学里面 有时候你的实验

其实就是去收集文献 对文献 对数据进行分析 那在有一些领域 不是只能做文献收集

不是只做文献收集 你有时候还需要建模型 训练模型 还要跑模型 今天AI agent

有办法跑模型吗 你可以看看前几天 安德鲍卡帕试出的 Auto Research 你可以拿一个LLM

让他自动帮你训练模型 他做的事情就是 横轴是每一次的实验 他大概五分钟做一次实验 你可以想一下

就好像他让他的agent 五分钟心跳一次 太过分了 我只让小金十五分钟心跳一次 他五分钟心跳一次

实在是太卷了 好 然后重轴呢 就是模型训练出来的模型的表现 那这个数值越低呢

代表结果越好 那这边的每一个点 蓝色绿色的点跟灰色的点 代表某一个模型的表现 那绿色的点代表说

被记录下来 结果比较好的模型 那在这整个过程中 没有人类的介入 就是叫一个AI agent

去训练模型 他先训练第一版的模型 然后看看 结果不够好 想想看

要改什么样的地方 改什么样 training script 哪边需要修改 再训练第二个模型

再训练第三个 一直训练下去 模型的结果 他自己训练出来的结果 就越来越好

这个过程中 不需要人类的介入 由AI来自主的 把一个模型训练好 那至于AI

到底可以做到 什么样的程度 你可以在我们的 作业二体验一下 我们的作业二

其实跟上学期 机器学习 导论的其中一个作业 是一模一样的 只是这一次

不是由人来完成 是由你操控 AI agent 来完成这个作业 你来看看

今天的语言模型 有没有能力 训练出自己的模型 好 那有人可能想说

刚才讲说 AI agent 可以整理文献 可以写文章 也可以做实验


\keyslide{../bilibili_video/slides_p3/slide_0009.jpg}{Slide 9}

\textbf{幻灯片引用链接:}
\begin{itemize}
  \item \href{https://github.com/karpathy/autoresearch}{https://github.com/karpathy/autoresearch}
\end{itemize}


\keyslide{../bilibili_video/slides_p3/slide_0010.jpg}{Slide 10}

\textbf{幻灯片引用链接:}
\begin{itemize}
  \item \href{https://arxiv.org/abs/2409.04109}{https://arxiv.org/abs/2409.04109}
\end{itemize}


\keyslide{../bilibili_video/slides_p3/slide_0011.jpg}{Slide 11}

\textbf{幻灯片引用链接:}
\begin{itemize}
  \item \href{https://arxiv.org/abs/2409.04109}{https://arxiv.org/abs/2409.04109}
\end{itemize}


\keyslide{../bilibili_video/slides_p3/slide_0012.jpg}{Slide 12}

\textbf{幻灯片引用链接:}
\begin{itemize}
  \item \href{https://arxiv.org/abs/2409.04109}{https://arxiv.org/abs/2409.04109}
\end{itemize}
\begin{itemize}
  \item \href{https://arxiv.org/a0s/2506.2C803}{https://arxiv.org/a0s/2506.2C803}
\end{itemize}


\keyslide{../bilibili_video/slides_p3/slide_0013.jpg}{Slide 13}

但是问题的发响 寻找问题 总是应该由人类来做吧 所以有人就写了一篇文章 在这篇文章是

可LLN 能不能够产生 新颖的研究的idea 这个其实不是很新的文章 这个是24年的文章

是古时候的文章 那个时候 LLN就展现了非常强的 产生研究idea的想法 他就让那个Language Model

去Pump Language Model 让他产生一大堆研究的想法 他让Language Model 对过去的论文 做一下RIG

然后产生一大堆研究的想法 接下来 他去找真正的人类 也产生研究的想法 接下来比较一下

AI跟人类 谁产生的研究想法 比较好 他就再另外找了一群人 来评价这些研究的想法

最后的结果是 他分出了几个不同的指标 来比较人类跟agent的差异 比如说Novelty 比如说excitement

比如说feasibility 比如说effectiveness 还有overall的指标 那这边有三个bro 第一个是人类

黄色的代表人类 浅蓝色的代表AI 你会发现在多数时候 人类其实是输给AI的 人类唯一赢过AI的

是feasibility 就是人类在想出来的idea的可行性上 比AI还要高 但如果是要讲novelty AI想的

居然由专家来评断 觉得AI想的研究题目 是比较有创新的 当然这个研究 如果你要批评的

你可以批评说 他们找来想研究题目的人不够强 其实他们在录文的附录里面 有告诉你说 他们找了什么样的人

来产生这些研究的题目 他还记录了这些人的H index 所以这些人也不是完全的麻瓜 他们其实都是领域里面的学者 如果没记错的话

他们是在国际会议 直接拉人来想这些题目的 但我觉得这边有可能的问题是 你在国际会议随便有一个人来说 你可不可以给我讲一个研究题目

你可能不会告诉他 你真正最好的研究题目 你想说这个人是要干嘛 想要脱我的研究题目吗 给你一个次等的

也许这边征求到的研究题目不是最好的研究题目 也说不定 如果你要批评的话 你可以这样讲 也许人类在这边还没有竭尽全力

这篇论文是有一个续作的 在一年之后 同样的团队做了一个续作 这个续作是 他们把人类跟AI提的这些点子

真的在找人去真的做成论文 每一篇就是做成大概四页左右的论文 然后再找另外一群人来审查这一些 由人类的idea还有AI的idea所产生出来的论文 这边一个有趣的现象

也许是让人类松一口气的现象是 当AI的点子真的被实作以后 看起来就没那么厉害了 这边灰色的现代表说 在实作之后

这一些点子 它的评分的改变 所以你会发现说 AI的点子 虽然在实作之前

人们觉得比人类还novel 但是真的实作之后 它的novel的程度 它的心颖的程度 就比不上人类了

所以整体而言的分数 AI想的那些题目 在人类真的实作以后 觉得就没有那么厉害了 因为AI有时候想的想法

是表面上看起来很厉害 堆砌很多新颖的词汇 让你觉得好像像是那么一回事 真的去执行的时候才发现 不太能够执行

做不太起来 所以AI agent 它产生出来的题目 其实最终是没有人类好的 也许这可以让人类松口气

但是不要忘了 那是2025年的时候 AI的能力不断的与日俱增 今天你可能觉得 GPT 5.4很厉害

或觉得CALCO 4.6很厉害 但是不要忘了 他们是2026年 现在这个时间点的未来 往后看最差的模型

以后我们只会看到 更厉害的模型而已 但至少在2025年 当第二篇论文论事的时候 那个时候

AI想的题目 实际上并没有人类想的题目好 那AI它当然也可以审查论文 今天你投稿到一个国际会议 就会有一群reviewer

一群审查委员 来对这篇论文进行评价 最后决定论文是否被国际会议接受 今天这一些reviewer 有没有可能


\keyslide{../bilibili_video/slides_p3/slide_0014.jpg}{Slide 14}

其实就是一个AI呢 其实在今年的Triple AI 这是一个跟AI有关的国际会议 AI是正式进入了审查流程 在Triple AI 2026里面


\keyslide{../bilibili_video/slides_p3/slide_0015.jpg}{Slide 15}

\textbf{幻灯片引用链接:}
\begin{itemize}
  \item \href{https://agents4science}{https://agents4science}
\end{itemize}

每篇文章 不只有三个人类的reviewer 还有一个AI的reviewer 他就是AI 他名字告诉你

我就是AI 然后呢 Meta Reviewer 也是一个人类的Meta Reviewer 一个AI的Meta Reviewer

不过跟人类不一样的地方是 这些AI呢 他们不打分数 他们只给意见 不过他们的意见

是人类最终做决定的时候 可以参考的 你可能想说 三个人类 一个AI

好像人类还比较多 但是你不知道 那些人类背后有多少 其实也是AI agent 我今年呢

Triple AI 我有担任这个area chair 在我负责的那些文章里面 就有一个review 他的第一句话是

Sure I can help you write this review 这样子 你了解他的意思吗 我反手就举报他了

举报这个reviewer 我在想说 这个会不会是一个测验 就是大会想要测验 area chair

有没有认真的在看所有的review 所以其中一些人类的review 虽然表面上有一个名字 但是他其实是AI 所以我必须要反手举报他

才能够展现说 我其实是真的有看这些review的 总之 你不知道有多少人类 背后其实是AI agent

那就我的立场而言 我其实并没有特别反对 用AI来辅助审查论文 我反对的是 拿不够好的AI来审查论文

比如说 我担任某个期刊的editor 然后呢 有人就给我 有一篇文章

有人给我一个很奇怪的review 那篇文章明明是一个分析的文章里面 没有propose任何方法 但是review的那个结果一直说 这篇文章propose了一个什么样的方法

根本就是牛头不对马嘴 然后呢 我就把那个review退回去 跟那个reviewer说 你写的东西牛头不对马嘴

比如说 你列的第一点 你说这篇文章提出来一个什么想法 但他其实没有提任何的想法 然后

那个reviewer居然过几分钟以后 就上传了新的review 我想说 哇 这个背后大概是个open cloud

但是 但是糟糕的地方是 他真的就只改了 他的第一点而已 我本来是想说举一个例子

因为后面都是错的 我举个例子告诉你 第一点都是错的 你后面要不要剪裁 但他用的语言模型太笨了

居然只改了第一点而已 感觉就是把我的instruction 丢给语言模型 让它产生一个新的review结果而已 我没有办法接受的

并不是用AI review 而是没有用一个最好的模型 来review这些文章 那有人可能听到AI review 就会觉得非常的反感

觉得说 哎呀 这个review审查 当然应该是要人类完成吗 怎么可以由AI代劳

但是你再想想review背后真正的意义 review的意义是什么 review的意义是 找出一篇文章的问题 让这篇文章变得更好

如果今天AI相较于人 他更能看出文章的问题 那为什么不让他做呢 如果人类根本就做得很差 那还不如由AI来得到更好的review结果

这样也许对论文的作者 还更有帮助 事实上小金在他成为一个YouTuber之前 他有一个工作 就是帮实验室的同学看论文

就是实验室的同学 在赶论文的时候 你可以把你的论文直接寄给小金 他会回复给你 他对于这篇论文的建议

但是后来 因为我们在赶一个 我们领域的国际会议 叫Interspeech 它是三月初解稿

论文解稿之后 小金就太闲了 所以就让他去变成一个YouTuber 本来他会每十分钟的收信一次 看看有没有人寄论文给他

如果有人寄论文给他 他就会给予回馈 那你可能会想说 这个跟学生直接把一篇paper 给ChairGPT做review

其实学生蛮常这么做的 今天我蛮常发现说 会有同学 直接把他写好的论文 丢给一个语言模型

然后看看语言模型 给他什么样的评价 然后他会根据语言模型的建议 来进行修改 把自己的文章做得更好

但对我来说 其实我觉得是一件好事 那小金的review 跟直接拿语言模型的review 有什么不同呢

那当然是因为 我在后面还是写了不同的指示 让我觉得小金的review 可以比其他人做得更好 举例来说

今天其他的模型 在review的时候 往往是批评多于建议 那小金在review的时候 小金review完一篇文章以后

他会把那个文章 他会把review的结果 也forward给我 所以我知道小金review的结果 长得什么样子

然后我一开始觉得 哇 这个review太过harsh了 太多的批评 太少的建议

对一个指导教授而言 假设站在指导教授的角色 你批评完以后 一定要有建议 不只说这篇文章很烂

你要说到底要怎么改 所以我就跟小金说 以后有批评 后面都要附上建议 后来他看起来的

review的感觉就好很多 或者是 在很接近deadline的时候 这个时候你还跟学生讲说 我觉得这个pand paper

应该要补一个实验 学生都会生气 很接近 比如说deadline还差24个小时 也许真正要做的就是

把文章的内容写好 而不是补做新的实验 所以我跟小金说 你在review的时候 要参考这个论文的deadline

要想说这个实验 如果要补实验 到底是不是能做的 如果今天离deadline很近 那你就只给

这个论文论事的建议就好了 所以小金就会照着做 如果离deadline很近 他就会给你论文的建议 而且如果离deadline越近

我就叫他给学生 多一点的鼓励 他就告诉你说 你论文写得非常很棒 我觉得这篇论文非常有希望

大家看了都得到情绪的价值 那能不能够改论文呢 我其实有把小金改论文的 一些原则 叫他做成一个影片

就是他上传的第二支影片 大家可以看看 对小金而言 他知不知道一篇好的论文 应该要长什么样子

那既然AI可以写论文 AI也可以审查论文 这就形成了一个闭环 一个国际会议 不需要人类

由AI来写论文投稿 由AI来审查论文决定是否接受 最后就产生一些高质量的被接受的论文 然后这国际会议就可以 可能以一个月一次的形式继续办下去

然后人类的技术就自动的进步了 中间不需要人类的介入 那确实有一群Stanford的研究人员 就做了类似的尝试 这个conference叫做

这个AI agent for science 就是这个conference AI必须要是第一作者 论文主要的贡献者 必须是AI


\keyslide{../bilibili_video/slides_p3/slide_0016.jpg}{Slide 16}

然后由AI来审查论文 这个会议的接受率 其实蛮低的 有247篇的投稿 最后只接受了48篇

接受率少小于20\% 跟多数的顶会一样 所以要被接受 是没有那么容易的 然后每篇论文

都会有三个AI reviewer给的分数 但最后他们还是找了人类 来给予一个最终的评价 那这些论文到底做得如何呢 在这群作者呢

后来有写了一篇paper summarize了整个AI conference的结果 AI投稿 AI审稿的结果 那在这个投稿的时候呢

每一个投稿者 都必须要indicate说 现在这篇文章 有多少AI的介入 那一般你在投稿

到人类的国际会议的时候 也有很多国际会议 会这样要求你 那通常都说 我只让AI帮我润稿而已

算其实有可能 做了更多的事情 在这篇paper里面 在这个国际会议里面 你必须要展示说

这篇文章 是由AI自主完成的 人类的介入越少越好 我是不知道我多少人虚报 就人类介入很多

但是其实谎称人类没有什么介入 然后他们把人类的介入 分成ABCD四个等级 第一代表说 有95\%以上

是由AI完成的 他们把产生一篇论文的工作 分成四个项目 包括这个点子 谁想的

实验 谁设计的 最后谁分析资料 最后谁写文章 那你会发现说

在这篇paper里面 有很大一部分 这四个项目 都是几乎由AI 自己完成的

但有趣的地方来了 这是所有paper的平均 如果我们看那些 被接受的paper 你会发现

被接受的paper 跟所有的paper比起来 在点子发想 跟实验设计上面 有非常明显的差距

今天被接受的那些paper 都是人类在点子发想 介入比较多的 还有人类在实验设计 介入比较多的

所以看起来 要产生一篇 比较好的文章 人类在初期的介入 还是需要的

那至于资料的分析 还有论文的写作 看起来AI 是有办法 独立完成的

那也有 在这篇论文里面 那这群作者 也提供了 投稿者的一些想法

比如说有一个投稿者 就说 对他而言 这些AI的问题 就是没办法

想出新颖的点子 没办法 产生有创造力的点子 多数时候 他们产生的点子


\keyslide{../bilibili_video/slides_p3/slide_0017.jpg}{Slide 17}

\textbf{幻灯片引用链接:}
\begin{itemize}
  \item \href{https://arxiv.org/abs/251}{https://arxiv.org/abs/251}
\end{itemize}

都没有那么新 都比较像是 把已有的东西 做一下重排 做一下组合而已

所以看起来 部分人类的介入 仍然是需要的 至少今天 以今天这一些

AI agent的能力 由人类来想出 比较好的问题 人类来想出 比较好的研究的方向


\keyslide{../bilibili_video/slides_p3/slide_0018.jpg}{Slide 18}

仍然是必要的 所以今天的状况 是这一些代理 这些agent 他可以自己完成任务

但是往往需要 人类来决定 来引导 他要完成 什么样的任务

由人类来告诉他 什么样的任务 才是真正重要的 我其实很好奇 今天这一些

AI agent 在教学上 可以做到 什么样的程度 其实我完全不怀疑

这些AI agent 是有教学的能力的 至少 他有做教学 投影片的能力

你可以看看 小金做的那些 教学影片 都是由他自主完成的 那并不是由我

协助他完成的 他做完以后 我也不会改 他做出来怎么样 就怎么样

这样子 所以他是由他 自主完成的 教学影片 他可以做到

那个程度 虽然跟人类 还是有一定程度的 差距 但是你想想看

他做的那个等级 人类如果不好好做的话 你是没有办法 讲的比AI还要更好的 但我想要知道

AI agent 在做教学影片上 可以做到 什么样的程度呢 那有一个比赛


\keyslide{../bilibili_video/slides_p3/slide_0019.jpg}{Slide 19}

这个是由台大 AI卓越中心 所办的比赛 这个比赛 就是招募

能够教学的AI 那这个比赛 会出题目 让AI 来做教学的影片

那如果你手上 有一个AI agent的话 你就告诉他 来参加这个比赛 看看他能做到

什么样的程度 好 那这个就是 今天上课的内容 那前两周呢

我们讲的是 比较科普的内容 讲AI agent 那下次上课呢 我们就要进入

语言模型 比较核心的部分 我们要看 语言模型内部 他是怎么做


\keyslide{../bilibili_video/slides_p3/slide_0020.jpg}{Slide 20}

\textbf{幻灯片引用链接:}
\begin{itemize}
  \item \href{https://youtu.be/8iFvVM7WUUs8?}{https://youtu.be/8iFvVM7WUUs8?}
\end{itemize}

Inference的 所以下次在上课之前 请先预习 至少先预习 我放在

头影片上的 这一部影片 这个是 机器学习 导论的第三讲

那你要看完第三讲 你才能够知道 下周要讲的内容 是什么 好

那以上就是 今天要跟大家 分享的内容


\end{document}